\cotefile{157-3}
\scanpage{130}
\mspage{79}

\uline{Remarque} La version duale \deleted{\unsure{du lieu|du lien}} (au sens de
Quillen) de l'axiome $\mathrm{W}(7\ \mathrm{bis})$ du carré cocartésien, i.e. l'axiome correspondant pour les
carrés cartésiens, n'est pas valable \deleted{\unsure{pour}} \unsure{$\mathcal W_\infty$}.
L'énoncé pertinent serait le suivant : soit

\diagram{\begin{tikzcd}[ampersand replacement=\&]
X \arrow[d,"f"'] \& X' \arrow[l,"q"'] \arrow[d,"f'"] \\
S \& S' \arrow[l,"p"']
\end{tikzcd}}

un carré cartésien, avec $f$ une \uline{$\mathcal W$-fibration}.
Alors $f\in\mathcal W\Longleftrightarrow f'\in\mathcal W$. Si \deleted{\unsure{Mais si on a seulement}
\unsure{soit} $f\in\mathcal W$} \unsure{:} \inserted{$\mathcal W$ satisfait} $\mathrm W(\unsure{4\,b|46})$\inserted{ (ce qui est le cas pour $\mathcal W_\infty$)}, alors $f\in\mathcal W$ implique
bien que ses fibres sont $\mathcal W$-asphériques, donc
que c'est une $\mathcal W$-équivalence universelle, i.e. $f'\in\mathcal W$
$\forall\,p:S'\unsure{-|\to}S$. Mais pour l'implication inverse
$f'\in\mathcal W\Rightarrow f\in\mathcal W$, il faut supposer \inserted{visiblement} que
\[
\pi_0(p):\pi_0(S')- \pi_0(S)
\]
est \uline{surjectif}. (P. ex. si $S'=\varnothing$, on a $X'=\varnothing$
donc $f'\in\mathcal W$, sans que ça implique $f\in\mathcal W$ !)

Le fait que dans le cas cocartésien de
$\mathrm W(7\ \mathrm{bis})$, $i_0\in\mathcal W\Rightarrow i\in\mathcal W$ doit être considéré
comme ``allant de soi'' -- autrement le carré
ne mérite pas, du point de vue $\mathcal W$-homotopique, le nom de ``cocartésien''. Par contre
l'implication inverse $\unsure{i|v}\in\mathcal W\unsure{\Longrightarrow|=}i_0\in\mathcal W$ m'apparaît comme un trait très particulier
du dérivateur $\mathrm{HOT}_{\mathcal W}$ (quand il le possède), et
notamment de $\mathrm{HOT}=\mathrm{HOT}_{\mathcal W_\infty}$. Le début

% archive sheet number: 130
